% English translation of chapter2.tex lines 116--292.
% Source: Methods of Algebra, Volume 2, commit
% 9a5803ff77dd3257484cb177f851a73770a59dd3 (CC BY 4.0).
% stable-unit-id: o014.aljabr2.chapter2.first-look-at-complexes

% segment-id: o014.li2.chapter2.first-look-at-complexes.h001
\section{A First Look at Complexes}\label{sec:cohomology}
\hypertarget{o014.aljabr2.chapter2.first-look-at-complexes}{ }

% segment-id: o014.li2.chapter2.first-look-at-complexes.p001
The starting point of this section is a pair of morphisms in an abelian
category satisfying $gf=0$, namely
$X \xrightarrow{f} Y \xrightarrow{g} Z$. The composite
$\Ker(g) \hookrightarrow Y \twoheadrightarrow \Coker(f)$ gives a morphism
$\Ker(g) \to \Coker(f)$; what interests us is its epi--mono factorization.
There are two mutually dual points of view.

% segment-id: o014.li2.chapter2.first-look-at-complexes.le001
\begin{lemma}\label{prop:image-to-kernel}
	Let $f:X\to Y$ and $g:Y\to Z$ be morphisms in an abelian category
	$\mathcal{A}$ such that $gf=0$.
	% segment-id: o014.li2.chapter2.first-look-at-complexes.l001
	\begin{enumerate}[(i)]
		% segment-id: o014.li2.chapter2.first-look-at-complexes.i001
		\item There are unique morphisms $\varphi$ and $\psi$ that make the
		following diagrams commute:
		% segment-id: o014.li2.chapter2.first-look-at-complexes.q001
		\begin{equation}\label{eqn:cplx-transition}
		% segment-id: o014.li2.chapter2.first-look-at-complexes.g001
		\begin{tikzcd}
			X \arrow[r, "f"] \arrow[twoheadrightarrow, d] & Y \\
			\Coim(f) \arrow[r, "\varphi"'] & \Ker(g) \arrow[hookrightarrow, u]
		\end{tikzcd} \qquad
		% segment-id: o014.li2.chapter2.first-look-at-complexes.g002
		\begin{tikzcd}
			Z & Y \arrow[twoheadrightarrow, d] \arrow[l, "g"'] \\
			\Image(g) \arrow[hookrightarrow, u] & \Coker(f) \arrow[l, "\psi"]
		\end{tikzcd}
		\end{equation}
		Moreover, $\varphi$ is a monomorphism and $\psi$ is an epimorphism.

		% segment-id: o014.li2.chapter2.first-look-at-complexes.i002
		\item The morphism $\Ker(g)\to\Coker(f)$ has the two epi--mono
		factorizations displayed in the following diagram; they are connected by
		a unique isomorphism $\Coker(\varphi)\simeq\Ker(\psi)$:
		% segment-id: o014.li2.chapter2.first-look-at-complexes.g003
		\[\begin{tikzcd}[row sep=tiny]
			& \Ker(\psi) \arrow[hookrightarrow, rd, "\text{canonical}"] \arrow[dd, "\sim" sloped] & \\
			\Ker(g) \arrow[twoheadrightarrow, ru] \arrow[twoheadrightarrow, rd, "\text{canonical}"'] & & \Coker(f) . \\
			& \Coker(\varphi) \arrow[hookrightarrow, ru] &
		\end{tikzcd}\]
	\end{enumerate}
\end{lemma}

% segment-id: o014.li2.chapter2.first-look-at-complexes.pf001
\begin{proof}
	First consider (i). The two diagrams are dual to one another
	(Remark~\ref{rem:im-coim-duality} and
	Proposition~\ref{prop:abelian-cat-duality}): the diagram on the right is
	equivalent to considering in $\mathcal{A}^{\opp}$ the sequence
	$Z \xrightarrow{g^{\opp}} Y \xrightarrow{f^{\opp}} X$. Thus, for
	\eqref{eqn:cplx-transition}, it suffices to consider $\varphi$.

	The condition $gf=0$ says precisely that $f$ factors through
	$\Ker(g)\hookrightarrow Y$. Applying
	Lemma~\ref{prop:Im-Coim-fac} on coimages to this factorization gives
	$\varphi$. Next, because $X\to\Coim(f)$ is an epimorphism, the
	factorization of $f$ as
	$X \twoheadrightarrow \Coim(f) \xrightarrow{\overline{f}} Y$ is unique.
	Here $\overline{f}$ may be taken either as the composite
	$\Coim(f) \xrightarrow{\varphi} \Ker(g) \hookrightarrow Y$ or as the
	composite $\Coim(f) \rightiso \Image(f) \hookrightarrow Y$. The latter
	composite is a monomorphism, and hence so is $\varphi$.

	To obtain (ii), we show that there is a unique monomorphism
	$\Coker(\varphi)\hookrightarrow\Coker(f)$
	(or epimorphism $\Ker(g)\twoheadrightarrow\Ker(\psi)$) that makes the
	following diagram commute:
	\[
	% segment-id: o014.li2.chapter2.first-look-at-complexes.g004
	\begin{tikzcd}[column sep=small]
		X \arrow[r, "f"] \arrow[twoheadrightarrow, d] & Y \arrow[twoheadrightarrow, r] & \Coker(f) \\
		\Coim(f) \arrow[hookrightarrow, r, "\varphi"'] & \Ker(g) \arrow[twoheadrightarrow, r, "\text{canonical}"' inner sep=0.6em] \arrow[hookrightarrow, u] & \Coker(\varphi) \arrow[hookrightarrow, u]
	\end{tikzcd} \quad
	% segment-id: o014.li2.chapter2.first-look-at-complexes.g005
	\begin{tikzcd}[column sep=small]
		Z & Y \arrow[twoheadrightarrow, d] \arrow[l, "g"'] & \Ker(g) \arrow[twoheadrightarrow, d] \arrow[hookrightarrow, l] \\
		\Image(g) \arrow[hookrightarrow, u] & \Coker(f) \arrow[twoheadrightarrow, l, "\psi"] & \Ker(\psi) \arrow[hookrightarrow, l, "\text{canonical}" inner sep=0.6em]
	\end{tikzcd}
	\]
	Once this is proved, the two epi--mono factorizations in (ii) can be read
	directly from the outer frames of the diagrams. The isomorphism along the
	middle path is then uniquely determined by
	Proposition~\ref{prop:epi-mono-uniqueness}.

	It remains to construct the diagrams above. Since the two diagrams are
	dual, it suffices to discuss the one on the left.
	Lemma~\ref{prop:mono-epi-ker-coker}~(iii) shows that
	$\Coker(\varphi)$ can also be identified with the cokernel of the composite
	$X\to\Ker(g)$. The desired morphism therefore comes from the functoriality
	of cokernels and is characterized by
	\eqref{eqn:Ker-Coker-functoriality}. Why is this morphism a monomorphism?
	Proposition~\ref{prop:Ker-Coker-composite} expresses the cokernel of a
	composite morphism as the pushout of the original cokernel; explicitly, the
	diagram is
	% segment-id: o014.li2.chapter2.first-look-at-complexes.g006
	\[\begin{tikzcd}
		\Ker(g) \arrow[hookrightarrow, r] \arrow[twoheadrightarrow, d] & Y \arrow[twoheadrightarrow, d] \\
		\Coker{[X \to \Ker(g)]} \arrow[r] & \Coker(f) \arrow[phantom, lu, "\boxplus" description]
	\end{tikzcd} \quad \text{(pushout diagram)}\]
	with all morphisms canonical. Proposition
	\ref{prop:Abel-cat-pull-push} then shows that the morphism in the second row
	is a monomorphism.
\end{proof}

% segment-id: o014.li2.chapter2.first-look-at-complexes.p002
In the situation of Lemma~\ref{prop:image-to-kernel}, write
\index[sym1]{HXYZ@$\Hm\left[X \to Y \to Z \right]$}
% segment-id: o014.li2.chapter2.first-look-at-complexes.q002
\begin{equation}\label{eqn:homology-3}\begin{aligned}
	\Hm\left[ X \xrightarrow{f} Y \xrightarrow{g} Z \right] & := \Coker\left[ \Image(f) \xrightarrow{\varphi} \Ker(g) \right] \\
	& \simeq \Ker\left[ \Coker(f) \xrightarrow{\psi} \Image(g) \right] ;
\end{aligned}\end{equation}
this object can also be characterized as the middle term in the epi--mono
factorization of $\Ker(g)\to\Coker(f)$. As expected,
\eqref{eqn:homology-3} is functorial.

% segment-id: o014.li2.chapter2.first-look-at-complexes.pr001
\begin{proposition}\label{prop:homology-functoriality}
	Given a commutative diagram in an abelian category
	% segment-id: o014.li2.chapter2.first-look-at-complexes.g007
	\[\begin{tikzcd}
		X \arrow[r, "f"] \arrow[d, "a"'] & Y \arrow[r, "g"] \arrow[d, "b"] & Z \arrow[d, "c"] \\
		X' \arrow[r, "{f'}"'] & Y' \arrow[r, "{g'}"'] & Z'
	\end{tikzcd} \qquad g'f' = 0, \quad gf = 0, \]
	there is a unique morphism
	$\Phi: \Hm[X \to Y \to Z] \to \Hm[X' \to Y' \to Z']$ that makes the
	following diagram commute:
	% segment-id: o014.li2.chapter2.first-look-at-complexes.g008
	\[\begin{tikzcd}
		\Ker(g) \arrow[twoheadrightarrow, r]  \arrow[d, "\text{induced by $(b,c)$}"'] & {\Hm[X \to Y \to Z]} \arrow[hookrightarrow, r] \arrow[d, "\Phi"] & \Coker(f) \arrow[d, "\text{induced by $(a,b)$}"] \\
		\Ker(g') \arrow[twoheadrightarrow, r] & {\Hm[X' \to Y' \to Z']} \arrow[hookrightarrow, r] & \Coker(f')
	\end{tikzcd}\]
	The two outer vertical arrows come from the functoriality of kernels and
	cokernels, while the horizontal arrows are those of
	Lemma~\ref{prop:image-to-kernel}.
\end{proposition}

% segment-id: o014.li2.chapter2.first-look-at-complexes.pf002
\begin{proof}
	As in the argument of Lemma~\ref{prop:epi-mono-morphism}, if $\Phi$ exists
	then it is unique, and it suffices to check that $\Phi$ makes the right-hand
	square commute. Here is the construction. Since
	$\Image(g)=\Ker[Z\to\Coker(g)]$ and
	$\Image(g')=\Ker[Z'\to\Coker(g')]$, the functoriality of kernels gives a
	unique morphism $\Image(g)\to\Image(g')$ that makes the right half of the
	following diagram commute:
	% segment-id: o014.li2.chapter2.first-look-at-complexes.g009
	\[\begin{tikzcd}
		\Coker(f) \arrow[d] \arrow[twoheadrightarrow, r, "\psi"] & \Image(g) \arrow[d] \arrow[hookrightarrow, r] & Z \arrow[d, "c"] \\
		\Coker(f') \arrow[twoheadrightarrow, r, "{\psi'}"'] & \Image(g') \arrow[hookrightarrow, r] & Z' ;
	\end{tikzcd}\]
	The pattern of Lemma~\ref{prop:epi-mono-morphism} shows that the left half
	also commutes. The morphism $\Ker(\psi)\to\Ker(\psi')$ induced by it is the
	desired morphism.
\end{proof}

% segment-id: o014.li2.chapter2.first-look-at-complexes.pr002
\begin{proposition}\label{prop:homology-biproduct}
	Given morphisms satisfying $gf=0$ and $g'f'=0$, there is a canonical
	isomorphism
	% segment-id: o014.li2.chapter2.first-look-at-complexes.q003
	\begin{multline*}
		\Hm\left[ X \xrightarrow{f} Y \xrightarrow{g} Z \right] \oplus \Hm\left[ X' \xrightarrow{f'} Y' \xrightarrow{g'} Z' \right] \\
		\simeq \Hm\left[ X \oplus X' \xrightarrow{f \oplus f'} Y \oplus Y' \xrightarrow{g \oplus g'} Z \oplus Z' \right].
	\end{multline*}
\end{proposition}

% segment-id: o014.li2.chapter2.first-look-at-complexes.pf003
\begin{proof}
	Use the morphisms $\iota$, $p$, and so forth associated with biproducts, as
	in \eqref{eqn:biproduct-identities}. By the functoriality in
	Proposition~\ref{prop:homology-functoriality}, these morphisms also induce
	% segment-id: o014.li2.chapter2.first-look-at-complexes.g010
	\[\begin{tikzcd}[column sep=small]
		\Hm{\left[ X \to Y \to Z \right]} \arrow[shift left, r, "\iota"] & \Hm{\left[ X \oplus X' \to Y \oplus Y' \to Z \oplus Z' \right]} \arrow[shift left, l, "p"] \arrow[shift right, r, "p"'] & \Hm\left[ X' \to Y' \to Z' \right] \arrow[shift right, l, "\iota"'].
	\end{tikzcd}\]

	The characterization of the induced morphism $\Phi$ in
	Proposition~\ref{prop:homology-functoriality} also tells us that
	$(a,b,c)\mapsto\Phi$ is compatible with composition, preserves addition,
	sends $0$ to $0$, and sends $\identity$ to $\identity$. Hence the morphisms
	induced at the level of $\Hm[\cdots]$ also satisfy the biproduct identities
	\eqref{eqn:biproduct-identities}.
\end{proof}

% segment-id: o014.li2.chapter2.first-look-at-complexes.p003
Notice that the preceding statement need not hold for infinite coproducts or
infinite products in a general abelian category.
% segment-id: o014.li2.chapter2.first-look-at-complexes.d001
\begin{definition}[Complexes]\label{def:complex}
	\index{complex}
	Let $\mathcal{A}$ be an additive category. Consider a sequence of morphisms
	in $\mathcal{A}$
	% segment-id: o014.li2.chapter2.first-look-at-complexes.q004
	\[ \cdots \to X^n \xrightarrow{d^n} X^{n+1} \xrightarrow{d^{n+1}} X^{n+2} \to \cdots. \]
	The sequence may have finitely many terms, extend infinitely in one or both
	directions, or form a cycle. If $d^{n+1}d^n=0$ for every $n$, the data
	$(X^n,d^n)_n$ are called a \emph{complex}; the customary notation is
	$(X^\bullet,d^\bullet)$, $X^\bullet$, or $X$. The term $X^n$ is
	conventionally called the degree-$n$ term of the complex $X$.
\end{definition}

% segment-id: o014.li2.chapter2.first-look-at-complexes.p004
This chapter is chiefly concerned with abelian categories, where the
cohomology and exactness of complexes can be studied.

% segment-id: o014.li2.chapter2.first-look-at-complexes.d002
\begin{definition}[Cohomology and exact sequences]\label{def:cohomology}
	\index{cohomology}
	\index{exact sequence}
	\index{complex!acyclic}
	\index{complex!exact}
	\index[sym1]{Hn@$\Hm^n$}
	Let $\mathcal{A}$ be an abelian category.
	% segment-id: o014.li2.chapter2.first-look-at-complexes.l002
	\begin{itemize}
		% segment-id: o014.li2.chapter2.first-look-at-complexes.i003
		\item For a complex $X$, in accordance with \eqref{eqn:homology-3}, its
		\emph{cohomology} at any non-end term $X^n$ is defined by
		% segment-id: o014.li2.chapter2.first-look-at-complexes.q005
		\[ \Hm^n(X) = \Hm^n(X^\bullet, d^\bullet) := \Hm\left[ X^{n-1} \xrightarrow{d^{n-1}} X^n \xrightarrow{d^n} X^{n+1}  \right]. \]
		% segment-id: o014.li2.chapter2.first-look-at-complexes.i004
		\item If $\Hm^n(X)=0$, the complex $X$ is said to be exact at $X^n$;
		this is equivalent to $\Image(d^{n-1})=\Ker(d^n)$. A complex that is
		exact at every term is called an \emph{exact sequence}; an exact complex
		is also called \emph{acyclic}.
	\end{itemize}
\end{definition}

% segment-id: o014.li2.chapter2.first-look-at-complexes.p005
The properties of cohomology are one of the main themes of this book and will
be studied further in \S\ref{sec:Abel-cplx}.

% segment-id: o014.li2.chapter2.first-look-at-complexes.r001
\begin{remark}[Homology of chain complexes]\label{rem:cochain-vs-chain}
	\index{complex!cochain}
	\index{complex!chain}
	\index{homology}
	\index[sym1]{Hnn@$\Hm_n$}
	In Definition~\ref{def:complex}, choose decreasing indices and use
	subscripts, namely
	% segment-id: o014.li2.chapter2.first-look-at-complexes.q006
	\[ \cdots \to X_n \xrightarrow{d_n} X_{n-1} \to \cdots, \quad d_{n-1} d_n = 0, \]
	and all the preceding definitions apply unchanged. By conventions originating
	in topology, the version with increasing indices $(X^\bullet,d^\bullet)$ is
	usually called a cochain complex, while the version with decreasing indices
	$(X_\bullet,d_\bullet)$ is called a chain complex. For a chain complex $X$
	in an abelian category, the \emph{homology} object at $X_n$ is defined as the
	following object of $\mathcal{A}$:
	% segment-id: o014.li2.chapter2.first-look-at-complexes.q007
	\[ \Hm_n(X) = \Hm_n(X_\bullet, d_\bullet) := \Hm\left[ X_{n+1} \xrightarrow{d_{n+1}} X_n \xrightarrow{d_n} X_{n-1} \right]. \]

	In algebra, complexes (that is, cochain complexes) and chain complexes differ
	only in notation, and passing between them is straightforward. For an
	additive category $\mathcal{A}$, one way to remain entirely within the same
	category is to reflect the indices by setting
	% segment-id: o014.li2.chapter2.first-look-at-complexes.q008
	\[ X^n := X_{-n}, \quad d^n := d_{-n}. \]
	Another way is to reverse the arrows, which involves the opposite category.
	If a complex in $\mathcal{A}$
	% segment-id: o014.li2.chapter2.first-look-at-complexes.q009
	\[ \cdots \to X^n \xrightarrow{d^n} X^{n+1} \to \cdots \]
	is viewed in $\mathcal{A}^{\opp}$, it becomes the chain complex
	% segment-id: o014.li2.chapter2.first-look-at-complexes.q010
	\begin{gather*}
		\cdots \leftarrow X_n \xleftarrow{d_{n+1}} X_{n+1} \leftarrow \cdots, \\
		X_n := X^n, \quad d_{n+1} := d^{n, \opp}.
	\end{gather*}

	If $\mathcal{A}$ is an abelian category and the chain complex
	$(X_\bullet,d_\bullet)$ in $\mathcal{A}^{\opp}$ is defined by reversing
	the arrows, there is a canonical isomorphism
	$\Hm^n(X^\bullet)\simeq\Hm_n(X_\bullet)$. In particular, $X^\bullet$ is
	exact if and only if $X_\bullet$ is exact. To see this, it suffices to note
	that $\Hm[X \xrightarrow{f} Y \xrightarrow{g} Z]$ can be characterized as
	the middle term in the epi--mono factorization of
	$\Ker(g)\to\Coker(f)$. Viewed in $\mathcal{A}^{\opp}$, the same object is
	also the middle term in the epi--mono factorization of
	$\Ker(f^{\opp})\to\Coker(g^{\opp})$.
\end{remark}

% segment-id: o014.li2.chapter2.first-look-at-complexes.p006
Exactness of a complex can be checked term by term. The following standard
forms will henceforth be used without further explanation.

% segment-id: o014.li2.chapter2.first-look-at-complexes.l003
\begin{description}
	% segment-id: o014.li2.chapter2.first-look-at-complexes.i005
	\item[Monomorphism] The complex $0\to X'\xrightarrow{f}X$ is exact if and
	only if $\Ker(f)=0$, that is, if and only if $f$ is a monomorphism
	(Lemma~\ref{prop:mono-epi-ker-coker}). This is also equivalent to
	$X'\rightiso\Coim(f)$ (Lemma~\ref{prop:epi-image}).

	% segment-id: o014.li2.chapter2.first-look-at-complexes.i006
	\item[Epimorphism] Dually, the complex
	$X\xrightarrow{g}X''\to0$ is exact if and only if $g$ is an epimorphism,
	and this is also equivalent to $\Image(g)\rightiso X''$.

	% segment-id: o014.li2.chapter2.first-look-at-complexes.i007
	\item[Kernel] Consider the complex
	$0\to X'\xrightarrow{f}X\xrightarrow{g}X''$. In accordance with
	\eqref{eqn:cplx-transition}, factor $f$ through
	$X'\twoheadrightarrow\Image(f)\hookrightarrow\Ker(g)$. The claim is that
	$0\to X'\to X\to X''$ is exact if and only if $X'\to\Ker(g)$ is an
	isomorphism. In other words, up to isomorphism, an exact sequence of this
	form is always $0\to\Ker(g)\to X\xrightarrow{g}X''$ and is given by the
	kernel of a morphism.

	The argument is easy. Consider \eqref{eqn:cplx-transition}. The complex is
	exact at $X$ if and only if $\Image(f)\rightiso\Ker(g)$, while exactness at
	$X'$ is equivalent to $X'\rightiso\Image(f)$. Now apply
	Proposition~\ref{prop:epi-mono-isomorphism} to
	$X'\twoheadrightarrow\Image(f)\hookrightarrow\Ker(g)$.

	% segment-id: o014.li2.chapter2.first-look-at-complexes.i008
	\item[Cokernel] Dually, the complex
	$X'\xrightarrow{f}X\xrightarrow{g}X''\to0$ induces a morphism
	$\Coker(f)\to X''$ through \eqref{eqn:cplx-transition}. This complex is
	exact if and only if $\Coker(f)\rightiso X''$. Thus, up to isomorphism, an
	exact sequence of this form arises from a cokernel.

	% segment-id: o014.li2.chapter2.first-look-at-complexes.i009
	\item[Short exact sequence] An exact sequence of the form
	$0\to X'\to X\to X''\to0$ is called a \emph{short exact sequence}. The two
	preceding items together give two equivalent, mutually dual
	characterizations of exactness:
	% segment-id: o014.li2.chapter2.first-look-at-complexes.l004
	\begin{compactitem}
		% segment-id: o014.li2.chapter2.first-look-at-complexes.i010
		\item $X'\to X$ is a monomorphism and
		$X''=\Coker[X'\hookrightarrow X]$;
		% segment-id: o014.li2.chapter2.first-look-at-complexes.i011
		\item $X\to X''$ is an epimorphism and
		$X'=\Ker[X\twoheadrightarrow X'']$.
	\end{compactitem}
	\index{short exact sequence}
\end{description}

% segment-id: o014.li2.chapter2.first-look-at-complexes.p007
As an application, every morphism $f:X\to Y$ gives short exact sequences
% segment-id: o014.li2.chapter2.first-look-at-complexes.q011
\[ 0 \to \Ker(f) \to X \to \Image(f) \to 0, \quad 0 \to \Image(f) \to Y \to \Coker(f) \to 0, \]
which can be spliced into the exact sequence
$0\to\Ker(f)\to X\xrightarrow{f}Y\to\Coker(f)\to0$.

% segment-id: o014.li2.chapter2.first-look-at-complexes.p008
Next, suppose that $f$ can be embedded in an exact sequence of the form
% segment-id: o014.li2.chapter2.first-look-at-complexes.q012
\[ L' \xrightarrow{\lambda} L \to X \xrightarrow{f} Y \to R \xrightarrow{\rho} R'' \]
with both $\lambda$ and $\rho$ isomorphisms. Exactness then forces both
$L\to X$ and $Y\to R$ to be zero, so $\Ker(f)=0$ and $\Image(f)=Y$; in other
words, $f$ is an isomorphism. This is a commonly used technique for testing
whether a morphism is an isomorphism.

% segment-id: o014.li2.chapter2.first-look-at-complexes.p009
The biproducts considered in \S\ref{sec:Ker-Coker} give a very simple class
of short exact sequences; Proposition \ref{prop:split-ses} will characterize
them.

% segment-id: o014.li2.chapter2.first-look-at-complexes.pr003
\begin{proposition}\label{prop:biproduct-ses}
	Let $X_1,X_2$ be objects of an abelian category $\mathcal{A}$. Then
	$0\to X_1\xrightarrow{\iota_1}X_1\oplus X_2\xrightarrow{p_2}X_2\to0$
	is a short exact sequence.
\end{proposition}

% segment-id: o014.li2.chapter2.first-look-at-complexes.pf004
\begin{proof}
	Take the short exact sequences
	$0\to X_1\xrightarrow{\identity}X_1\to0\to0$ and
	$0\to0\to X_2\xrightarrow{\identity}X_2\to0$.
	Proposition~\ref{prop:homology-biproduct} shows that their direct sum is
	still exact.
\end{proof}
